\section{Limitations}

The scientific evaluation is limited to symbolic, score-level SATB harmonization. The product pipeline can render score-derived audio for playback validation, but audio synthesis quality is not evaluated as a scientific endpoint in this paper. The work does not address accompaniment tracks, drums, or popular-music arrangements. MIDI pitch numbers are used only as internal symbolic pitch tokens. The dataset is the Bach chorale collection bundled with music21, so the reported evidence should not be generalized to all choral repertories without additional corpora.

The harmonic annotations are approximate automatic labels. Roman numerals, chord qualities, seventh-chord flags, dominant-function flags, and cadence labels are derived with music21 and local pitch-class heuristics. Ambiguous sonorities are marked as \texttt{UNKNOWN}; they are not converted into artificial ground truth. Consequently, seventh-resolution and cadence metrics should be interpreted as conservative diagnostics rather than definitive musicological analysis.

The primary RTX rows and logged ablations are available, but the logged vanilla Transformer and Transformer-without-constraints rows refer to the same no-constraints checkpoint. They should therefore not be interpreted as independent baselines. The experiments use one deterministic split and one logged seed, so they do not establish variance across seeds or corpora. Expert-evaluation ratings are still pending, so claims about stylistic preference, singability, and pedagogical value remain unconfirmed. Finally, the symbolic decoder is a local repair layer, not an exhaustive search procedure; generation-level validity, long-sequence handling, and conflicting harmony-and-counterpoint objectives remain open limitations.
